\begin{chapterpage}{Inferensi untuk data numerik}
  \chaptertitle{Inferensi untuk data numerik}
  \label{inferenceForNumericalData}
  \label{ch_inference_for_means}
  \chaptersection{oneSampleMeansWithTDistribution}
  \chaptersection{pairedData}
  \chaptersection{differenceOfTwoMeans}
  \chaptersection{PowerForDifferenceOfTwoMeans}
  \chaptersection{anovaAndRegrWithCategoricalVariables}
\end{chapterpage}
\renewcommand{\chapterfolder}{ch_inference_for_means}


\chapterintro{Bab~\ref{ch_foundations_for_inf}
  memperkenalkan kerangka inferensi statistik yang didasarkan
  pada interval kepercayaan dan hipotesis menggunakan
  distribusi normal untuk proporsi sampel.
  Dalam bab ini, kita akan menjumpai beberapa estimasi titik baru
  dan dua distribusi baru.
  Dalam setiap kasus, gagasan inferensinya tetap sama:
  tentukan estimasi titik atau statistik uji yang berguna,
  identifikasi distribusi yang sesuai untuk estimasi titik
  atau statistik uji tersebut, lalu
  terapkan gagasan inferensi.}



%__________________
\section[Rata-rata satu sampel dengan distribusi $t$]
    {Rata-rata satu sampel dengan distribusi
        $\pmb{\MakeLowercase{t}}$}
\label{oneSampleMeansWithTDistribution}

\noindent%
Sebagaimana kita dapat memodelkan perilaku
proporsi sampel $\hat{p}$ menggunakan distribusi normal,
rata-rata sampel $\bar{x}$ juga dapat dimodelkan menggunakan
distribusi normal ketika syarat-syarat tertentu terpenuhi.
\index{estimasi titik!rata-rata tunggal}
Namun, kita akan segera mempelajari bahwa distribusi baru,
yang disebut distribusi $t$,
cenderung lebih berguna ketika bekerja dengan rata-rata sampel.
Pertama-tama kita akan mempelajari distribusi baru ini,
kemudian menggunakannya untuk menyusun interval kepercayaan
dan melakukan uji hipotesis untuk rata-rata.


\subsection[Distribusi $\bar{x}$]
    {Distribusi sampling dari $\pmb{\bar{x}}$}

Rata-rata sampel cenderung mengikuti
distribusi normal yang berpusat pada rata-rata populasi,~$\mu$,
ketika syarat-syarat tertentu terpenuhi.
Selain itu, kita dapat menghitung galat baku untuk rata-rata
sampel menggunakan simpangan baku populasi $\sigma$
dan ukuran sampel $n$.

\begin{onebox}{Teorema Limit Pusat untuk rata-rata sampel}
  Ketika kita mengumpulkan sampel cukup besar yang terdiri atas
  $n$~observasi yang saling independen dari populasi dengan
  rata-rata $\mu$ dan simpangan baku $\sigma$,
  distribusi sampling $\bar{x}$ akan mendekati
  distribusi normal dengan
  \begin{align*}
  &\text{Rata-rata}=\mu
  &&\text{Galat Baku }(SE) = \frac{\sigma}{\sqrt{n}}
  \end{align*}
\end{onebox}

\noindent%
Sebelum mendalami interval kepercayaan dan uji
hipotesis menggunakan $\bar{x}$, pertama-tama kita perlu membahas dua topik:
\begin{itemize}
\item
    Ketika kita memodelkan $\hat{p}$ menggunakan distribusi normal,
    syarat-syarat tertentu harus terpenuhi.
    Syarat-syarat untuk bekerja dengan $\bar{x}$
    sedikit lebih rumit, dan kita akan menggunakan
    Bagian~\ref{x_bar_conditions} untuk membahas
    cara memeriksa syarat-syarat inferensi.
\item
    Galat baku bergantung pada simpangan
    baku populasi, $\sigma$.
    Namun, kita jarang mengetahui $\sigma$, sehingga
    kita harus mengestimasinya.
    Karena estimasi ini sendiri tidak sempurna,
    kita menggunakan distribusi baru yang disebut
    distribusi $t$\index{distribusi t@distribusi $t$}
    untuk mengatasi masalah ini, yang kita bahas dalam
%    While we can use the plug-in principle,
%    using the sample standard deviation $s$ in place of $\sigma$,
%    this is not quite enough to resolve the issue entirely.
%    and .
%    We'll cover this topic in
    Bagian~\ref{introducingTheTDistribution}.
\end{itemize}


\subsection[Mengevaluasi dua syarat yang diperlukan untuk
    memodelkan $\bar{x}$]
  {Mengevaluasi dua syarat yang diperlukan untuk
      memodelkan $\pmb{\bar{x}}$}
\label{x_bar_conditions}

\noindent%
Dua syarat diperlukan untuk menerapkan
Teorema Limit Pusat\index{Teorema Limit Pusat}
pada rata-rata sampel~$\bar{x}$:
\begin{description}
\item[Independensi.]
    Observasi sampel harus saling independen.
    Cara paling umum untuk memenuhi syarat ini adalah
    mengambil sampel acak sederhana dari
    populasi.
    Jika data berasal dari proses acak,
    yang serupa dengan pelemparan sebuah dadu,
    hal ini juga akan memenuhi syarat independensi.
\item[Normalitas.]
    Ketika sampelnya kecil,
    kita juga mensyaratkan bahwa observasi sampel
    berasal dari populasi yang berdistribusi normal.
    Syarat ini semakin dapat dilonggarkan
    seiring bertambahnya ukuran sampel.
    Syarat ini memang bersifat samar,
    sehingga sulit dievaluasi;
    karena itu, selanjutnya kita memperkenalkan dua pedoman praktis
    agar pemeriksaan syarat ini menjadi lebih mudah.
\end{description}
%%Before we get to the sample size consideration, let's
%%consider a special case of the normal distribution
%%where any sample size is sufficient.
%
%%There is also a special case of the Central Limit Theorem
%%for when the data come from a nearly normal distribution.
%%In this case the sample mean will be nearly normal
%%regardless of sample size.
%
%\begin{onebox}{Special case of the Central Limit Theorem
%    for normally distributed data}
%  The sampling distribution of $\bar{x}$ is nearly normal when
%  the sample observations are independent and come from a nearly
%  normal distribution.
%  This is true for any sample size.
%\end{onebox}
%
%%For population distributions that are not normal,
%%the sample mean $\bar{x}$ will still look normal if the sample
%%size is large enough.
%%To check what is \emph{large enough}, we ask two questions:
%%\begin{itemize}
%%\item
%%    Is the sample show evident skew or outliers?
%%    If so, then if t
%%\end{itemize}
%
%In practice, the population never exactly follows
%a normal distribution,
%and the more ``non-normal'' a population
%distribution, the larger the required sample size required for
%$\bar{x}$ to be reasonably modeled using a normal distribution.
%The rough rule of thumb is, if you don't see any clear outliers
%and we don't have reason to believe particularly extreme outliers
%are present in population, then this condition is satisfied.

\begin{onebox}{Pedoman praktis:
    cara melakukan pemeriksaan normalitas}
  Tidak ada cara sempurna untuk memeriksa syarat normalitas,
  jadi sebagai gantinya kita menggunakan dua pedoman praktis: %,
%  one for small samples ($n < 30$)
%  and another for large samples ($n \geq 30$):
  \begin{description}
  \setlength{\itemsep}{0mm}
  \item[$\mathbf{n < 30}$:]
      Jika ukuran sampel $n$ kurang dari 30
      dan tidak terdapat pencilan yang jelas dalam data,
      biasanya kita mengasumsikan data berasal dari
      distribusi yang mendekati normal agar
      syarat tersebut terpenuhi.
  \item[$\mathbf{n \geq 30}$:]
      Jika ukuran sampel $n$ setidaknya 30
      dan tidak terdapat pencilan yang \emph{sangat ekstrem},
      biasanya kita mengasumsikan distribusi sampling
      $\bar{x}$ mendekati normal, sekalipun distribusi
      observasi individual yang mendasarinya tidak normal.
  \end{description}
\end{onebox}

Dalam mata kuliah statistika pengantar ini, Anda tidak diharapkan
mengembangkan penilaian yang sempurna mengenai syarat normalitas.
Namun, Anda diharapkan mampu menangani
kasus-kasus yang jelas berdasarkan pedoman praktis tersebut.\footnote{Pedoman
  yang lebih cermat akan mempertimbangkan pelonggaran lebih lanjut
  atas pemeriksaan \emph{pencilan sangat ekstrem} ketika
  ukuran sampel sangat besar.
  Namun, pembahasan lebih lanjut akan kita sisakan untuk mata kuliah berikutnya.}

\begin{examplewrap}
\begin{nexample}{Perhatikan dua grafik berikut
    yang berasal dari sampel acak sederhana dari
    populasi yang berbeda.
    Ukuran sampelnya adalah $n_1 = 15$ dan $n_2 = 50$.
    \begin{center}
    \Figure[Dua histogram ditampilkan, satu untuk "Observasi Sampel 1" dan satu lagi untuk "Observasi Sampel 2". Histogram Observasi Sampel 1 memuat nilai dari 0 hingga 7 dengan lebar bin 1, sehingga terdapat 7 bin dengan frekuensi 2, 1, 4, 3, 2, 0, dan 3. Histogram Observasi Sampel 2 memuat nilai dari 0 hingga 22, dengan lebar bin 1. Sebagian besar data berada dekat nol, dengan separuh observasi berada dalam bin dari 0 hingga 1. Hanya ada satu bin tak nol di atas 5; tingginya tampak 1, yaitu bin dari 21 hingga 22.]{0.85}{outliers_and_ss_condition}
    \end{center}
    Apakah syarat independensi dan normalitas terpenuhi
    dalam setiap kasus?}
  \label{outliers_and_ss_condition_ex}%
  Setiap sampel merupakan sampel acak sederhana dari
  populasinya masing-masing, sehingga syarat independensi
  terpenuhi.
  Selanjutnya, mari kita periksa syarat normalitas untuk
  setiap sampel menggunakan pedoman praktis.
  
  Sampel pertama memiliki kurang dari 30 observasi,
  sehingga kita memperhatikan setiap pencilan yang jelas.
  Tidak ada pencilan semacam itu; meskipun terdapat celah kecil dalam
  histogram antara 5 dan~6, celah ini kecil dan
  20\% observasi dalam sampel kecil ini
  terwakili oleh batang paling kanan pada histogram,
  sehingga observasi itu nyaris tidak dapat disebut pencilan yang jelas.
  Karena tidak ada pencilan yang jelas, syarat normalitas
  secara wajar~terpenuhi.

  Sampel kedua berukuran lebih besar dari 30 dan
  mencakup pencilan yang tampaknya kira-kira 5 kali
  lebih jauh dari pusat distribusi dibandingkan
  observasi terjauh berikutnya.
  Ini merupakan contoh pencilan yang sangat ekstrem,
  sehingga syarat normalitas tidak akan terpenuhi.
\end{nexample}
\end{examplewrap}
